% !TeX root = ../AMaN_SUN.tex
% 全项目通用可调设置；不限于数学排版。

% 正文段间距；目录等紧凑区域可在局部覆盖。
\newlength{\BodyParagraphSpacing}
\setlength{\BodyParagraphSpacing}{3ex plus 0.3ex minus 0.2ex}

% 公式局部强调。
% Highlight：行内/行间公式分别设置横向、纵向内边距；浓度取值为 0--100。
\newlength{\MathHighlightInlineHorizontalPadding}
\setlength{\MathHighlightInlineHorizontalPadding}{3.2pt}
\newlength{\MathHighlightInlineVerticalPadding}
\setlength{\MathHighlightInlineVerticalPadding}{2.2pt}

\newlength{\MathHighlightDisplayHorizontalPadding}
\setlength{\MathHighlightDisplayHorizontalPadding}{4pt}
\newlength{\MathHighlightDisplayVerticalPadding}
\setlength{\MathHighlightDisplayVerticalPadding}{3.5pt}
\newcommand{\MathHighlightTint}{16}

% Frame：行内/行间公式分别设置横向、纵向内边距，并共用边框粗细。
\newlength{\MathFrameInlineHorizontalPadding}
\setlength{\MathFrameInlineHorizontalPadding}{1.2pt}
\newlength{\MathFrameInlineVerticalPadding}
\setlength{\MathFrameInlineVerticalPadding}{1pt}

\newlength{\MathFrameDisplayHorizontalPadding}
\setlength{\MathFrameDisplayHorizontalPadding}{1.6pt}
\newlength{\MathFrameDisplayVerticalPadding}
\setlength{\MathFrameDisplayVerticalPadding}{2pt}
\newlength{\MathFrameRuleWidth}
\setlength{\MathFrameRuleWidth}{0.65pt}

% 全局等号颜色。
\definecolor{MathEqualsColor}{HTML}{909090}
